\section{Cài đặt system call \texttt{sysinfo}}
\label{sec:sysinfo}

\subsection{Yêu cầu và ý tưởng}

\cmd{sysinfo(struct sysinfo *info)} thu thập thông tin hệ thống đang chạy:
\begin{itemize}
  \item Các file đang mở trong toàn hệ thống, tức các ô \cmd{ftable.file[]}
        có \cmd{ref > 0}: số lượng và với mỗi file là type, mode, reference count,
        danh sách process tham chiếu.
  \item Các process có \cmd{state != UNUSED} trong \cmd{proc[]}: số lượng và với
        mỗi process là pid, state, pid của cha, số lần được lập lịch, danh sách
        file đang mở.
\end{itemize}

Danh sách ``process tham chiếu tới file'' được tính ngược: process nào có một
ô \cmd{ofile[]} trỏ tới \cmd{struct file} đó. Để người dùng đối chiếu được hai
danh sách, mỗi file mang chỉ số \cmd{idx} trong \cmd{ftable}, và mỗi fd của
process mang chỉ số \cmd{fileidx} của ô nó trỏ tới.

\subsection{Cấu trúc dữ liệu (\texttt{kernel/sysinfo.h})}

\begin{table}[H]
\centering
\small
\begin{tabularx}{\textwidth}{>{\ttfamily}l X}
\toprule
\multicolumn{2}{l}{\cmd{struct fileinfo} -- một file đang mở} \\
\midrule
idx & Vị trí trong \cmd{ftable.file[]} \\
type, readable, writable, ref & Lấy trực tiếp từ \cmd{struct file} \\
inum, itype, major, off & Số inode, loại inode (file/dir/device), thiết bị,
  offset; $-1$ nếu không áp dụng (ví dụ pipe) \\
nfd & Tổng số fd (của mọi process) trỏ vào file này \\
npid, pids[NPROC] & Danh sách pid các process tham chiếu (mỗi pid một lần) \\
\midrule
\multicolumn{2}{l}{\cmd{struct procinfo} -- một process khác \cmd{UNUSED}} \\
\midrule
pid, state, ppid, name & Định danh, trạng thái, pid của cha (0 nếu không có), tên \\
nsched & Số lần được scheduler chọn chạy (field mới trong \cmd{struct proc}) \\
nofile, fds[], fileidx[] & Các fd đang mở và chỉ số ô \cmd{ftable} tương ứng \\
\midrule
\multicolumn{2}{l}{\cmd{struct sysinfo}} \\
\midrule
nfile, files[NFILE] & Số file đang mở và mảng \cmd{fileinfo} \\
nproc, procs[NPROC] & Số process và mảng \cmd{procinfo} \\
\bottomrule
\end{tabularx}
\caption{Các cấu trúc trả về cho user}
\end{table}

\subsection{Các thay đổi trong mã nguồn}

\begin{table}[H]
\centering
\small
\begin{tabularx}{\textwidth}{>{\ttfamily}l X}
\toprule
\textrm{\textbf{File}} & \textbf{Thay đổi} \\
\midrule
kernel/sysinfo.h & (mới) Định nghĩa \cmd{fileinfo}, \cmd{procinfo}, \cmd{sysinfo}. \\
kernel/proc.h & Thêm field \cmd{nsched} vào \cmd{struct proc}. \\
kernel/proc.c & \cmd{allocproc}: \cmd{nsched = 0}. \cmd{scheduler}: \cmd{nsched++}
  trước \cmd{swtch}. Thêm \cmd{procfilerefs()} và \cmd{procsysinfo()}. \\
kernel/file.c & Thêm \cmd{fileindex()} và \cmd{filesysinfo()}. \\
kernel/sysproc.c & Thêm \cmd{sys\_sysinfo()}. \\
kernel/syscall.h, syscall.c & \cmd{SYS\_sysinfo = 23}; thêm vào bảng \cmd{syscalls[]}
  và \cmd{syscallnames[]} (để lệnh \cmd{trace} của Lab0 in được tên). \\
kernel/defs.h & Prototype của 4 hàm mới. \\
user/user.h, user/usys.pl & Khai báo \cmd{struct sysinfo;},
  \cmd{int sysinfo(struct sysinfo *);} và stub \cmd{entry("sysinfo")}. \\
user/sysinfotest.c & (mới) Chương trình test. \\
user/sysinfo.c & (mới) Lệnh in trạng thái hệ thống. \\
Makefile & Thêm \cmd{\_sysinfotest}, \cmd{\_sysinfo} vào \cmd{UPROGS}. \\
\bottomrule
\end{tabularx}
\caption{Danh sách file thay đổi}
\end{table}

\subsection{Luồng xử lý trong kernel}

\begin{enumerate}
  \item \cmd{sys\_sysinfo()} lấy địa chỉ user \cmd{addr} bằng \cmd{argaddr(0)}.
  \item Gọi \cmd{filesysinfo(addr + offsetof(files))}. Với mỗi ô của
        \cmd{ftable.file[]}:
        lấy \cmd{ftable.lock}; bỏ qua nếu \cmd{ref == 0}; chép type, mode, ref,
        thông tin inode vào biến cục bộ \cmd{fileinfo}; gọi
        \cmd{procfilerefs(f)} để duyệt \cmd{proc[]} và ghi lại pid các process
        có \cmd{ofile[]} trỏ tới \cmd{f}; nhả khóa; \cmd{copyout} phần tử đó
        vào \cmd{files[n]} của user.
  \item Gọi \cmd{procsysinfo(addr + offsetof(procs))}. Với mỗi \cmd{p} trong
        \cmd{proc[]}: lấy \cmd{wait\_lock} rồi \cmd{p->lock}; bỏ qua nếu
        \cmd{UNUSED}; chép pid, state, \cmd{parent->pid}, nsched, name và danh
        sách fd (kèm \cmd{fileindex()} của từng file) vào biến cục bộ
        \cmd{procinfo}; nhả khóa; \cmd{copyout} vào \cmd{procs[n]}.
  \item \cmd{copyout} hai giá trị \cmd{nfile}, \cmd{nproc} vào đầu struct và
        trả về 0. Nếu bất kỳ \cmd{copyout} nào thất bại (địa chỉ không hợp lệ,
        không ghi được hoặc vượt \cmd{MAXVA}) thì trả về $-1$.
\end{enumerate}

\subsection{Các vấn đề kỹ thuật và cách giải quyết}

\paragraph{Struct quá lớn so với stack.}
\cmd{struct sysinfo} có kích thước khoảng 40\,KB
($100 \times 300$\,B $+ 64 \times 164$\,B), trong khi kernel stack chỉ có 4\,KB.
Vì vậy kernel không dựng toàn bộ struct rồi \cmd{copyout} một lần, mà chỉ giữ
\emph{một} \cmd{fileinfo} hoặc \cmd{procinfo} trên stack và \cmd{copyout}
thẳng vào đúng vị trí trong struct của user
(\cmd{addr + offsetof(files) + n * sizeof(fileinfo)}). Ở phía user, user
stack chỉ có 8\,KB nên chương trình test khai báo struct là biến
\cmd{static} toàn cục.

\paragraph{Đồng bộ trên nhiều CPU.}
xv6 chạy trên 3 CPU, nên trong lúc \cmd{sysinfo} đọc, process khác có thể đang
\cmd{open}, \cmd{close}, \cmd{fork} hay \cmd{exit}. Nhóm dùng các khóa sẵn có
theo đúng vai trò của chúng:
\begin{itemize}
  \item \cmd{ftable.lock} được giữ trong suốt lúc đọc một ô \emph{và} dò
        \cmd{ofile[]}, nên \cmd{ref} và danh sách pid luôn nhất quán
        (\cmd{nfd} không bao giờ vượt \cmd{ref}).
  \item \cmd{p->lock} bảo vệ \cmd{state}, \cmd{pid}; \cmd{wait\_lock} bảo vệ
        \cmd{parent}.
  \item Thứ tự khóa: phần file là \cmd{ftable.lock} $\to$ \cmd{p->lock}; trong
        xv6 không có chỗ nào giữ \cmd{p->lock} rồi mới lấy \cmd{ftable.lock}
        (\cmd{exit} đóng file trước khi lấy \cmd{p->lock}), nên không gây
        deadlock. Phần process là \cmd{wait\_lock} $\to$ \cmd{p->lock}, cùng
        thứ tự với \cmd{wait()} và \cmd{exit()}.
  \item Mọi khóa được nhả trước khi \cmd{copyout}.
\end{itemize}

\paragraph{Đếm số lần lập lịch.}
\cmd{nsched} được tăng trong \cmd{scheduler()} ngay trước \cmd{swtch} vào
process, là nơi duy nhất một process được cấp CPU, và lúc đó đang giữ
\cmd{p->lock} (chi tiết ở mục~\ref{sec:sched}).

\subsection{Kiểm thử}
\label{sec:testing}

\cmd{user/sysinfotest.c} gọi \cmd{sysinfo()} trước và sau mỗi thao tác rồi so
sánh hai kết quả. Mỗi lần gọi còn được kiểm tra các bất biến chung: \cmd{ref > 0},
\cmd{nfd} $\le$ \cmd{ref}, process \cmd{ZOMBIE} không giữ file, mọi
\cmd{fileidx} của process đều có trong danh sách file và file đó liệt kê
ngược lại pid của process.

\begin{table}[H]
\centering
\small
\begin{tabularx}{\textwidth}{>{\ttfamily}l X}
\toprule
\textrm{\textbf{Test}} & \textbf{Kịch bản và kết quả mong đợi} \\
\midrule
basic & Thấy \cmd{init} (pid 1, ppid 0, SLEEPING), \cmd{sh} là cha của test,
  test ở trạng thái RUNNING; fd 0/1/2 cùng trỏ một file console
  (\cmd{FD\_DEVICE}, major 1, mode \cmd{rw}). \\
bad addresses & Địa chỉ 0 (trang text chỉ đọc), địa chỉ vượt \cmd{MAXVA},
  struct tràn qua cuối heap đều trả về $-1$; gọi đúng sau đó vẫn thành công. \\
open / close & \cmd{nfile} tăng/giảm 1; mode đúng với \cmd{O\_RDONLY},
  \cmd{O\_WRONLY}, \cmd{O\_RDWR}; mở thư mục có \cmd{itype} dir; offset là 10
  sau khi ghi 10 byte; hai lần \cmd{open} cùng file cho hai ô khác nhau cùng inode. \\
dup & Không tạo ô mới; 3 fd cùng \cmd{idx}, \cmd{ref = nfd = 3}; offset dùng
  chung; \cmd{close} giảm \cmd{ref}. \\
pipe & \cmd{nfile} tăng 2; một đầu chỉ đọc, một đầu chỉ ghi; không có inode. \\
fork / exit / wait & Con có \cmd{ppid} đúng, SLEEPING khi chờ pipe, cùng
  \cmd{idx} với cha, \cmd{ref} tăng (console tăng 3); con exit nhưng chưa
  wait thì là ZOMBIE, không còn file, \cmd{ref} giảm; sau \cmd{wait} con biến mất. \\
many processes & 10 process con: \cmd{nproc} tăng 10; đầu đọc pipe có
  \cmd{ref = npid = 11}; đầu ghi chỉ cha giữ. \\
scheduling counter & \cmd{nsched} tăng ít nhất 5 sau 5 lần \cmd{sleep(1)};
  process con mới có \cmd{nsched} nhỏ. \\
fd limit & Mở tới khi lỗi: \cmd{nofile} đúng bằng \cmd{NOFILE} = 16. \\
ftable full & 12 process con mở file tới khi lỗi: \cmd{nfile} đúng bằng
  \cmd{NFILE} = 100; sau khi con thoát, các ô được giải phóng. \\
\bottomrule
\end{tabularx}
\caption{Các nhóm test trong \texttt{sysinfotest}}
\end{table}

Kết quả: tất cả test đều đạt, chạy hai lần liên tiếp trong cùng một lần khởi
động vẫn đạt (không rò rỉ file hay process). \cmd{usertests -q} của xv6 cũng
đạt toàn bộ, cho thấy các thay đổi không ảnh hưởng tới phần còn lại của kernel.

\begin{figure}[H]
\begin{minted}{text}
$ sysinfotest
test basic:
  open files: 1
    idx type   mode ref inum itype off  pids
    0   device rw   9   23   dev   0    1 2 3
  processes: 3
    pid ppid state  nsched name        fd->ftable
    1   0    sleep  13     init        0->0 1->0 2->0
    2   1    sleep  12     sh          0->0 1->0 2->0
    3   2    run    23     sysinfotest 0->0 1->0 2->0
  OK
test bad addresses:
  OK
test open / close / modes:
  OK
test dup:
  OK
test pipe:
  OK
test fork / exit / wait:
  ...
  OK
test many processes:
  OK
test scheduling counter:
  nsched 68 -> 73
  OK
test per-process fd limit (NOFILE):
  OK
test system-wide file table full (NFILE):
  nfile=100 (NFILE=100)
  OK
sysinfotest: ALL TESTS PASSED
\end{minted}
\caption{Kết quả chạy \texttt{sysinfotest} trên QEMU}
\label{fig:test-result}
\end{figure}

\subsection{Chạy với chương trình thật}
\label{sec:demo}

Lệnh \cmd{sysinfo [-f | -p]} (\cmd{user/sysinfo.c}) gọi system call và in bảng
file (\cmd{-f}), bảng process (\cmd{-p}) hoặc cả hai. Kết hợp với shell, lệnh
này cho thấy trực quan các khái niệm ở phần tìm hiểu.

\paragraph{Pipeline.}
Với \cmd{cat README | grep the | sysinfo} (Hình~\ref{fig:demo-pipe}), shell tạo
hai pipe (ô 1--2 và 3--4). stdout của \cmd{cat} (fd 1) trỏ vào ô 2 (đầu ghi
pipe thứ nhất), stdin của \cmd{grep} trỏ vào ô 1 (đầu đọc), stdout của
\cmd{grep} trỏ vào ô 4, stdin của \cmd{sysinfo} trỏ vào ô 3. \cmd{cat} mở
README ở fd 3 (ô 5, đã đọc tới offset 2048). Console (ô 0) có
\cmd{ref = 16} vì được mọi process kế thừa qua \cmd{fork}. Các process
\cmd{sh} (pid 28, 30) là các bước fork trung gian của shell.

\begin{figure}[H]
\begin{minted}{text}
$ cat README | grep the | sysinfo
OPEN FILES: 6
IDX  TYPE    MODE  REF  INUM  ITYPE  OFF   PIDS
0    device  rw    16   23    dev    0     1 2 28 29 30 31 32
1    pipe    r-    2    -     -      -     30 31
2    pipe    -w    1    -     -      -     29
3    pipe    r-    1    -     -      -     32
4    pipe    -w    1    -     -      -     31
5    inode   r-    1    2     file   2048  29

PROCESSES: 7
PID  PPID  STATE  NSCHED  NAME     FD->FILE
1    0     sleep  13      init     0->0 1->0 2->0
2    1     sleep  15      sh       0->0 1->0 2->0
28   2     sleep  1       sh       0->0 1->0 2->0
29   28    sleep  8       cat      0->0 1->2 2->0 3->5
30   28    sleep  1       sh       0->1 1->0 2->0
31   30    run    9       grep     0->1 1->4 2->0
32   30    run    9       sysinfo  0->3 1->0 2->0
\end{minted}
\caption{\texttt{sysinfo} ở cuối một pipeline}
\label{fig:demo-pipe}
\end{figure}

\paragraph{Zombie và reparent.}
Chương trình \cmd{zombie} của xv6 fork một con rồi \cmd{sleep}; con thoát ngay.
Chạy \cmd{zombie \&; sysinfo -p} (Hình~\ref{fig:demo-zombie}): pid 36 ở trạng
thái ZOMBIE và không còn fd nào, chờ cha (pid 35) gọi \cmd{wait}. Pid 35 có
\cmd{ppid = 1} vì process shell chạy nền đã tạo nó đã thoát, nên nó được
\cmd{reparent} cho \cmd{init}.

\begin{figure}[H]
\begin{minted}{text}
$ zombie &; sysinfo -p
PROCESSES: 5
PID  PPID  STATE   NSCHED  NAME     FD->FILE
1    0     sleep   14      init     0->0 1->0 2->0
2    1     sleep   17      sh       0->0 1->0 2->0
33   2     run     2       sysinfo  0->0 1->0 2->0
36   35    zombie  1       zombie
35   1     sleep   6       zombie   0->0 1->0 2->0
\end{minted}
\caption{Process ZOMBIE và process được giao cho \texttt{init}}
\label{fig:demo-zombie}
\end{figure}

\paragraph{Chuyển hướng output.}
\cmd{sysinfo -f > out} rồi \cmd{cat out} (Hình~\ref{fig:demo-redirect}): file
\cmd{out} (inode 24) chiếm ô 1 với mode chỉ ghi, và chính là fd 1 của
\cmd{sysinfo}, đúng như phân tích ở bước 7 của mục~\ref{sec:trace-open}.
Lệnh \cmd{trace} của Lab0 với mask $2^{23}$ cũng in được tên syscall mới.

\begin{figure}[H]
\begin{minted}{text}
$ sysinfo -f > out
$ cat out
OPEN FILES: 2
IDX  TYPE    MODE  REF  INUM  ITYPE  OFF  PIDS
0    device  rw    8    23    dev    0    1 2 37
1    inode   -w    1    24    file   0    37
$ trace 8388608 sysinfo -p
syscall sysinfo -> 0
...
\end{minted}
\caption{Chuyển hướng output vào file và \texttt{trace} syscall \texttt{sysinfo}}
\label{fig:demo-redirect}
\end{figure}

\subsection{Khó khăn và hạn chế}

\begin{itemize}
  \item \textbf{Kích thước struct:} ban đầu dự định dựng \cmd{struct sysinfo}
        trong kernel rồi \cmd{copyout} một lần, nhưng 40\,KB vượt quá kernel stack
        4\,KB. Giải pháp là \cmd{copyout} từng phần tử.
  \item \textbf{\texttt{printf} của xv6 không hỗ trợ \texttt{\%c}:} bản đầu
        của bảng in ra ký tự rác ở cột mode, phải đổi sang in chuỗi.
  \item \textbf{Test bảng file đầy:} bản đầu dùng 8 process con nhưng chỉ đạt
        93 file, vì mỗi con đã có sẵn 5 fd nên chỉ mở thêm được 11 file. Đã tăng
        lên 12 process con.
  \item \textbf{Đọc \texttt{ip->type} không khóa:} trường này được bảo vệ bởi
        sleep-lock của inode, nhưng không thể ngủ khi đang giữ spinlock
        \cmd{ftable.lock}. Vì file đang mở (\cmd{ref > 0}) nên inode luôn hợp lệ
        và loại của nó không đổi, nên việc đọc là an toàn.
  \item \textbf{Ảnh chụp tại một thời điểm:} kết quả chỉ đúng tại lúc gọi; các
        process khác có thể thay đổi trạng thái ngay sau đó. Mỗi phần tử là nhất
        quán, nhưng giữa các phần tử với nhau thì không được bảo đảm
        (không khóa toàn bộ hệ thống trong suốt lời gọi).
  \item Các field \cmd{pids[]} và \cmd{fds[]} có kích thước cố định theo
        \cmd{NPROC} và \cmd{NOFILE}, nên struct khá lớn dù đa số phần tử để trống.
\end{itemize}
